\section{Модель распределённой системы}

В данной работе распределённая система рассматривается как совокупность
независимых вычислительных узлов, взаимодействующих посредством обмена
сообщениями по каналам связи и совместно реализующих единое функциональное
поведение. В соответствии с определением Э.
Таненбаума\cite{tanenbaum_distributed_systems}:

«Распределенная вычислительная система (РВС) – это набор соединенных каналами
связи независимых компьютеров, которые с точки зрения пользователя некоторого
программного обеспечения выглядят единым целым».

Такое представление соответствует рассматриваемой в работе системе
распределённых вычислений, в которой множество узлов совместно обрабатывает
задачи, обменивается служебной информацией и поддерживает согласованное
состояние. Каждый узел функционирует как самостоятельный вычислительный
компонент, однако итоговое поведение системы определяется не отдельными узлами,
а их согласованным взаимодействием.

В разрабатываемой системе каждый узел обладает собственным локальным
состоянием, выполняет предусмотренные алгоритмом действия и может отправлять
сообщения другим узлам. Общее состояние системы формируется совокупностью
состояний всех узлов, а также результатами их взаимодействия при передаче
команд, распределении ролей и обработке вычислительных задач. Внешние
пользователи или клиентские процессы могут направлять системе запросы,
инициирующие изменение её состояния и запуск вычислений.

Для дальнейшего анализа важны три группы предположений: свойства каналов связи,
временная модель системы и допустимые типы отказов. Именно эти предположения
определяют применимость алгоритмов консенсуса и границы корректного
функционирования распределённой системы.

\subsection{Модель каналов связи}

В распределённых системах обмен сообщениями не является абсолютно надёжным:
сообщения могут задерживаться, теряться, приходить повторно или наблюдаться в
различном порядке. В качестве базовой абстракции часто используется канал с
допустимыми потерями (fair-loss link) \cite{cachin11}. Для него характерны
следующие свойства:

\begin{itemize}
    \item если корректный процесс отправляет сообщение бесконечно много раз, то
корректный получатель в конечном итоге получит его бесконечное число раз;
    \item одно отправленное сообщение не может быть доставлено бесконечное
число раз;
    \item канал не создаёт новых сообщений.
\end{itemize}

Данная абстракция полезна как исходная модель ненадёжного сетевого
взаимодействия, однако её недостаточно для построения практических алгоритмов
консенсуса. Поэтому в распределённых протоколах обычно используются
дополнительные механизмы повышения надёжности: повторные передачи,
подтверждения доставки, идентификаторы сообщений и средства устранения
дубликатов.

С учётом этих механизмов вводится более сильная абстракция --- совершенный
канал связи (perfect link), обеспечивающий следующие гарантии \cite{cachin11}:

\begin{itemize}
    \item каждое сообщение, отправленное корректным процессом корректному
получателю, в конечном итоге будет доставлено;
    \item сообщения между двумя процессами доставляются без нарушения порядка;
    \item сообщение не доставляется более одного раза;
    \item канал не создаёт сообщений, которые не были отправлены.
\end{itemize}

При этом даже совершенный канал не устраняет всех проблем распределённого
взаимодействия. В реальной системе возможны временная недоступность связи,
длительные задержки и разделение сети, при котором часть узлов на некоторый
промежуток времени не может обмениваться сообщениями с остальными.
Следовательно, при построении модели распределённой системы необходимо
учитывать не только свойства отдельного канала, но и возможность нарушения
связности сети в целом.

В рамках настоящей работы каналы связи рассматриваются на уровне абстракции,
достаточном для анализа алгоритма консенсуса и поведения распределённой системы
вычислений. Существенным является не моделирование низкоуровневых сетевых
деталей, а фиксация того факта, что обмен сообщениями может сопровождаться
задержками, временной недоступностью и повторными передачами, что
непосредственно влияет на выбор лидера, репликацию состояния и согласованность
узлов системы.

\subsection{Консенсус в распределённой системе}

Центральной задачей распределённой системы, в которой несколько узлов совместно
поддерживают единое состояние, является достижение консенсуса. Под консенсусом
понимается согласование общего решения между множеством процессов, каждый из
которых действует независимо и обладает лишь частичной информацией о глобальном
состоянии системы.

Корректный алгоритм консенсуса должен обеспечивать выполнение следующих свойств
\cite{petrov}:

\begin{itemize}
    \item \textit{согласованность} --- все корректные процессы принимают одно и
то же решение;
    \item \textit{действительность} --- принятое значение должно быть
допустимым и происходить из множества предложенных значений;
    \item \textit{окончательность} --- процесс, принявший решение, не изменяет
его в дальнейшем, а алгоритм должен приводить систему к устойчивому результату.
\end{itemize}

Для алгоритмов консенсуса эти свойства имеют фундаментальное значение,
поскольку именно они обеспечивают согласованность реплик, корректность выбора
лидера и единый порядок применения команд. В настоящей работе консенсус
рассматривается как основа согласования состояния узлов распределённой системы
вычислений, обеспечивающая единый порядок применения команд, корректность
выбора лидера и устойчивость системы к частичным отказам.

\subsection{Временная модель и ограничения достижимости консенсуса}

Существенное влияние на возможность достижения консенсуса оказывает временная
модель распределённой системы. В полностью асинхронной системе отсутствуют
гарантированные верхние оценки времени доставки сообщений и времени выполнения
шагов процесса. В таких условиях узел не может определить, отказал ли другой
узел, либо его ответ лишь задерживается.

Фундаментальным результатом для таких систем является теорема
Фишера--Линча--Патерсона о невозможности детерминированного консенсуса в
полностью асинхронной среде даже при наличии всего одного отказа процесса
\cite{fischer85}. Этот результат показывает, что алгоритм, гарантирующий
достижение консенсуса за конечное время при любых асинхронных задержках, в
общем случае невозможен.

На практике распределённые системы обычно рассматриваются не как полностью
асинхронные, а как частично синхронные. Для такой модели предполагается, что
после некоторого неизвестного момента времени задержки сообщений и скорости
выполнения процессов оказываются ограниченными сверху, либо такие ограничения
выполняются в большинстве типичных сценариев \cite{cachin11}. Именно это
предположение лежит в основе большинства практически применимых невизантийских
алгоритмов консенсуса.

Следовательно, при анализе распределённой системы необходимо учитывать наличие
временного аспекта, однако не требуется вводить полностью синхронную модель с
жёстко фиксированными задержками. Для рассматриваемого класса систем достаточно
предположения о частичной синхронности, при котором алгоритм консенсуса
сохраняет работоспособность в типичных сценариях.

\subsection{Модели отказов}

Неотъемлемой характеристикой распределённой системы является модель отказов,
определяющая, какие нарушения поведения узлов и каналов связи допускаются в
рамках анализа. От выбора модели отказов зависит как проектирование алгоритма
консенсуса, так и архитектура самой распределённой системы.

В данной работе рассматриваются невизантийские модели отказов. Это означает,
что узлы системы либо работают в соответствии с алгоритмом, либо временно
перестают выполнять требуемые действия, но не пытаются произвольно искажать
данные и не действуют злонамеренно.

Основными типами отказов, существенными для настоящей работы, являются
следующие:

\begin{itemize}
    \item \textit{аварийное завершение} --- процесс прекращает выполнение шагов
алгоритма и перестаёт участвовать в обмене сообщениями;
    \item \textit{восстановление после сбоя} --- после отказа процесс может
быть перезапущен и синхронизирован с текущим состоянием системы, если алгоритм
предусматривает сохранение устойчивого состояния \cite{skeen83};
    \item \textit{пропуск действий} --- отдельные шаги алгоритма или результаты
этих шагов могут не наблюдаться другими узлами из-за потерь сообщений,
перегрузки или временной недоступности связи.
\end{itemize}

Византийские ошибки в рамках данной работы не рассматриваются. Это означает,
что не моделируются ситуации, в которых узел произвольно искажает сообщения,
распространяет противоречивую информацию разным участникам или преднамеренно
нарушает правила протокола.

Для невизантийских алгоритмов консенсуса основным средством повышения
отказоустойчивости является избыточность. За счёт репликации
состояния и принятия решений на основе большинства система способна
сохранять корректность при отказе части узлов \cite{christian91}.

\subsection{Принятые в работе допущения}

На основании рассмотренных свойств распределённых систем в данной работе
принимаются следующие допущения:

\begin{itemize}
    \item система состоит из конечного множества взаимодействующих узлов,
каждый из которых обладает локальным состоянием и может обмениваться
сообщениями с другими узлами;
    \item узлы совместно поддерживают согласованное состояние на основе
невизантийского алгоритма консенсуса;
    \item в системе допускаются задержки сообщений, временная недоступность
связи и разделение сети;
    \item временная модель считается частично синхронной;
    \item существенными типами отказов являются аварийное завершение, временная
недоступность узла и пропуск отдельных действий или сообщений;
    \item византийские отказы не рассматриваются;
    \item модель системы должна быть достаточна для описания поведения
выбранного алгоритма консенсуса и основных сценариев функционирования
распределённой системы вычислений.
\end{itemize}

Таким образом, используемая в работе модель распределённой системы
ориентирована не на максимально полное воспроизведение всех аспектов сетевой
среды, а на выделение тех свойств, которые существенны для анализа алгоритма
консенсуса и проектирования отказоустойчивой системы распределённых вычислений.
